%!TEX program = xelatex
\let\stop\empty

\documentclass[12pt,a4paper,fontset=none]{exam-zh}

% ==========================
% 卷面变量设置（只需改这里）
% ==========================
% ==========================
% 卷面变量设置（只需改这里）
% ==========================
\newcommand{\paperlevel}{初中一}
\newcommand{\papersubject}{数学 试卷一}
\newcommand{\schoolname}{金宝培元独立中学}
\newcommand{\examname}{年终考}
\newcommand{\classinfo}{J1(H / I/ F)}
\newcommand{\examdate}{30-10-2026}
\newcommand{\examtime}{13:00-14:15}
\newcommand{\questiontotal}{6}
\newcommand{\paperpercentage}{60\%}
\newcommand{\setterteacher}{蔡杰庆师}
\newcommand{\reviewteacher}{禤佩娟师}
\newcommand{\secondreviewteacher}{张惠雯师}
\newcommand{\auditteacher}{雷泳健师}



\usepackage{xeCJK}
\setCJKmainfont{NSimSun}
\newCJKfontfamily\heiti{NSimSun}
\usepackage{geometry}
\usepackage{array}
\usepackage{tikz}
\usepackage{amsmath}
\usepackage{wrapfig}
\usepackage{adjustbox}
\usepackage{enumitem}
\usepackage{fancyhdr}
\usetikzlibrary{quotes,angles}
\usetikzlibrary{calc}
\usepackage{scalerel}
\def\mcirc{\mathbin{\scalerel*{\circ}{j}}}
\def\msquare{\mathord{\scalerel*{\Box}{gX}}}

% \renewcommand{\choiceshook}{%
%     \setlength{\leftmargin}{18pt}%
% }
\geometry{margin=2cm}
\AtBeginDocument{\zihao{4}}


%region ----------- 调整 exam -----------
\examsetup{
    % problem/label = \arabic*., 
    % 设置大题题号格式 
        problem = {
        top-sep = 1em,      % 题头上方的间距（默认 0.25em）
        bottom-sep = 1em,   % 题号下方的间距（默认 0pt）
    },
    % 隐藏大题后的分值显示 
    problem/show-points = false,
    % 调整各 question 之间的垂直间距
    question = {
        top-sep = 1em,
        bottom-sep = 1em,
    },
    % 调整小题列表的缩进和间距
    list/step-name = {}, 
    solution = {
    show-solution = show-move, % 答案自动移动到卷末
    label-content = {参考答案},  % 答案区的标题文字
    show-qed = true            % 在解答末尾显示结束符号 ◻ 
  },
    page = {
    size = a4paper,           % 设定纸张大小 
    foot-content = { - ; - }  % 页脚显示为 -1-, -2- 格式 
  }
}
\setlist[enumerate,1]{
    label=(\alpha*), 
    itemsep=0.5em, 
    leftmargin=2em, 
    after=\hfill % 配合您之前要求的百分比靠右样式
}
%endregion

\newcommand{\qscoreoffset}{4em}
\newlist{questionparts}{enumerate}{1}
\setlist[questionparts]{
    label=(\alph*),
    itemsep=0.5em,
    leftmargin=2em,
    rightmargin=4em,
    before={\renewcommand{\qscoreoffset}{8em}},
}
\newcommand{\qscore}[2][\%]{%
    \leavevmode\unskip\nobreak\hfill
    \makebox[0pt][l]{\makebox[\qscoreoffset][r]{(#2#1)}}%
}

\newcommand{\makeexamheader}{%
%region 三栏标题区
% ==========================
\noindent
\begin{tabular*}{\textwidth}{@{\extracolsep{\fill}} p{0.3\textwidth} p{0.35\textwidth} p{0.3\textwidth}}

\raggedright
\textbf{\paperlevel}\\[0.2em]
\textbf{}{\papersubject}\\[0.4em]
姓名：\underline{\hspace{3cm}} &

\centering
\textbf{\schoolname}\\[0.2em]
\textbf{\examname}\\[0.4em]
班级：\classinfo &

\raggedright
日期：\examdate\\[0.2em]
时间：\examtime\\[0.4em]
得分：\underline{\hspace{2.5cm}} \\

\end{tabular*}

\begin{flushleft}
作答题：\questiontotal 题全做（\paperpercentage） 
\end{flushleft}
%endregion
}

\begin{document}
\makeexamheader


%region 范围：进位制与数制转换 | 考点：二进数与八进数互换（三位分组法）、二进数加法后转换 | 难度：易
\begin{question}[points = 8]
    请完成下列各题：

    \begin{enumerate}[label=(\alph*)]
        \item 将二进数 $1101011_2$ 转换为八进数；\qscore{3}

        \item 将八进数 $3746_8$ 转换为二进数；\qscore{3}

        \item 计算 $1011011_2 + 475_8$，答案以八进数表示。\qscore{4}
    \end{enumerate}
\end{question}

\begin{solution}
    \begin{enumerate}[label=(\alph*)]
        \item 将 $1101011_2$ 转换为八进数。 \\
        从右向左每三位分为一组，不足三位的在左侧补 $0$，再将每组转为八进数字。 \\
        \[
            1101011_2 = \underbrace{001}_{1}\ \underbrace{101}_{5}\ \underbrace{011}_{3}{}_{2}
        \]
        \[
            1101011_2 = 153_8
        \]

        \item 将 $3746_8$ 转换为二进数。 \\
        将每一位八进数字分别写成三位二进数，再依次连接，最高位的 $0$ 可省略。 \\
        \[
            3746_8 = \underbrace{011}_{3}\ \underbrace{111}_{7}\ \underbrace{100}_{4}\ \underbrace{110}_{6}{}_{2}
        \]
        \[
            3746_8 = 11111100110_2
        \]

        \item 求 $1011011_2 + 475_8$，并将和以八进数表示。 \\
        先把 $475_8$ 转为二进数，再按二进数逢二进一相加，最后将和从右向左每三位分组转为八进数。 \\
        \[
            475_8 = \underbrace{100}_{4}\ \underbrace{111}_{7}\ \underbrace{101}_{5}{}_{2} = 100111101_2
        \]
        \[
            \begin{array}{r}
                \phantom{+\ }001011011 \\
                +\ 100111101 \\
                \hline
                \phantom{+\ }110011000
            \end{array}
        \]
        \[
            1011011_2 + 475_8 = 110011000_2 = \underbrace{110}_{6}\ \underbrace{011}_{3}\ \underbrace{000}_{0}{}_{2} = 630_8
        \]
    \end{enumerate}
\end{solution}
%endregion

%region 范围：1.6.1 比例性质、1.6.2 连比 | 考点：比例的基本性质（含带分数）、连比的求法、连比求代数式的值 | 难度：易→中
\begin{question}[points = 10]
    请完成下列各题：

    \begin{enumerate}[label=(\alph*)]
        \item 已知 $2(x-1):4\dfrac{1}{2} = 5:3$，求 $x$ 的值；\qscore{3}

        \item 已知 $a:b = 5:6$，$b:c = 10:9$，求 $a:b:c$；\qscore{3}

        \item 根据(b)，若 $a+b+c = 164$，求 $2a+3b-4c$ 的值。\qscore{4}
    \end{enumerate}
\end{question}

\begin{solution}
    \begin{enumerate}[label=(\alph*)]
        \item 已知 $2(x-1):4\dfrac{1}{2} = 5:3$，求 $x$。 \\
        先把带分数化为假分数，再利用比例的基本性质：内项之积等于外项之积。 \\
        \begin{align*}
            4\dfrac{1}{2} &= \dfrac{9}{2} \\
            3\times 2(x-1) &= \dfrac{9}{2} \times 5
        \end{align*}
        \begin{align*}
            6(x-1) &= \dfrac{45}{2} \\
            x-1 &= \dfrac{15}{4}
        \end{align*}
        \begin{align*}
            x &= \dfrac{15}{4} + 1 = \dfrac{19}{4} = 4\dfrac{3}{4}
        \end{align*}

        \item 已知 $a:b = 5:6$，$b:c = 10:9$，求 $a:b:c$。 \\
        将两个比中 $b$ 所对应的数化为相同（取 $6$ 与 $10$ 的最小公倍数 $30$），再连写。 \\
        \begin{align*}
            a:b &= 5:6 = (5\times5):(6\times5) = 25:30 \\
            b:c &= 10:9 = (10\times3):(9\times3) = 30:27
        \end{align*}
        \begin{align*}
            a:b:c &= 25:30:27
        \end{align*}

        \item 已知 $a:b:c = 25:30:27$ 且 $a+b+c=164$，求 $2a+3b-4c$。 \\
        按连比设比例系数，先由和求出系数，再代入所求式。 \\
        设 $a = 25k$，$b = 30k$，$c = 27k$（$k>0$）。 \\
        \begin{align*}
            25k + 30k + 27k &= 164 \\
            82k &= 164 \\
            k &= 2
        \end{align*}
        \begin{align*}
            2a+3b-4c &= 2\times25k + 3\times30k - 4\times27k \\
            &= 50k + 90k - 108k
        \end{align*}
        \begin{align*}
            2a+3b-4c &= 32k = 32\times2 = 64
        \end{align*}
    \end{enumerate}
\end{solution}
%endregion

%region 范围：1.6.3 正比例、反比例、比例分配应用题 | 考点：正比例、反比例的应用 | 难度：中
\begin{question}[points = 10]
    请完成下列各题：

    \begin{enumerate}[label=(\alph*)]
        \item 某汽车行驶 $150$ km 耗油 $12$ 升。则该车行驶 $400$ km 需耗油多少升？\qscore{5}

        \item $20$ 名工人 $35$ 天可完成一项工程，每名工人工作效率相同。如果要提早 $10$ 天完成，需再增加多少名工人？\qscore{5}
    \end{enumerate}
\end{question}

\begin{solution}
    \begin{enumerate}[label=(\alph*)]
        \item 求行驶 $400$ km 的耗油量。 \\
        耗油量与路程成正比，故两者之比相等。 \\
        设行驶 $400$ km 耗油 $x$ 升。 \\
        \begin{align*}
            \dfrac{12}{150} &= \dfrac{x}{400}
        \end{align*}
        \begin{align*}
            150x &= 12 \times 400
        \end{align*}
        \begin{align*}
            x &= 32
        \end{align*}

        \item 求为提早 $10$ 天完成需再增加的工人数。 \\
        工人数与所需天数成反比，即工人数与天数之积（总工作量）不变。 \\
        设再增加 $y$ 名工人。 \\
        \begin{align*}
            \text{总工作量} &= 20 \times 35 = 700 \ \text{（人·天）}
        \end{align*}
        \begin{align*}
            \text{提早后所需天数} &= 35 - 10 = 25 \ \text{（天）}
        \end{align*}
        \begin{align*}
            (20 + y) \times 25 &= 700 \\
            20 + y &= 28
        \end{align*}
        \begin{align*}
            y &= 8
        \end{align*}
    \end{enumerate}
\end{solution}
%endregion



%region 范围：2.1.1 求值、2.1.2 代数式运算 | 考点：去括号与合并同类项、化简后代入求值、列代数式表示实际情境 | 难度：中
\begin{question}[points = 10]
    请完成下列各题：

    \begin{enumerate}[label=(\alph*)]
        \item 化简 $3(2x-y) - 2(x-3y) + (x+y)$；\qscore{4}

        \item 若$x = -3$，求$(3x^2 - 2x + 1) - 2(x^2 - 3x - 2)$的值。\qscore{6}

    \end{enumerate}
\end{question}

\begin{solution}
    \begin{enumerate}[label=(\alph*)]
        \item 化简 $3(2x-y) - 2(x-3y) + (x+y)$。 \\
        先去括号（括号前为负号时各项变号），再合并同类项。 \\
        \begin{align*}
            3(2x-y) - 2(x-3y) + (x+y) &= 6x - 3y - 2x + 6y + x + y
        \end{align*}
        \begin{align*}
            &= (6x - 2x + x) + (-3y + 6y + y) \\
            &= 5x + 4y
        \end{align*}

        \item 化简 $(3x^2 - 2x + 1) - 2(x^2 - 3x - 2)$ 并求 $x=-3$ 时的值。 \\
        先去括号合并同类项化简，再代入 $x=-3$。 \\
        \begin{align*}
            (3x^2 - 2x + 1) - 2(x^2 - 3x - 2) &= 3x^2 - 2x + 1 - 2x^2 + 6x + 4
        \end{align*}
        \begin{align*}
            &= (3x^2 - 2x^2) + (-2x + 6x) + (1 + 4) \\
            &= x^2 + 4x + 5
        \end{align*}
        \begin{align*}
            \text{当 } x = -3 \text{ 时，}\quad x^2 + 4x + 5 &= (-3)^2 + 4\times(-3) + 5
        \end{align*}
        \begin{align*}
            &= 9 - 12 + 5 \\
            &= 2
        \end{align*}

        \item 求应找回的钱数，并在 $x=8$，$y=4$ 时求值。 \\
        找回的钱数等于付出的 $100$ 元减去所买物品的总价。 \\
        \begin{align*}
            \text{总价} &= 3x + 5y \\
            \text{找回} &= 100 - (3x + 5y)
        \end{align*}
        \begin{align*}
            \text{当 } x = 8,\ y = 4 \text{ 时，}\quad 100 - (3\times8 + 5\times4) &= 100 - 44 \\
            &= 56 \ \text{元}
        \end{align*}
    \end{enumerate}
\end{solution}
%endregion

\newpage

%region 范围：2.5.1 一元一次方程及应用题 | 考点：含分数的一元一次方程（去分母、去括号）、年龄问题列方程 | 难度：中→难
\begin{question}[points = 10]
    请完成下列各题：

    \begin{enumerate}[label=(\alph*)]
        \item 解方程式$\dfrac{2x+1}{3} - \dfrac{5x-1}{6} = 1$；\qscore{5}

        \item $3$ 年前，父亲的年龄是儿子的 $7$ 倍。若儿子现在是 $10$ 岁，问父亲现在几岁？\qscore{5}
    \end{enumerate}
\end{question}

\begin{solution}
    \begin{enumerate}[label=(\alph*)]
        \item 解方程 $\dfrac{2x+1}{3} - \dfrac{5x-1}{6} = 1$。 \\
        方程两边同乘分母的最小公倍数 $6$ 去分母，再去括号（括号前为负号，各项变号）。 \\
        \begin{align*}
            6\times\dfrac{2x+1}{3} - 6\times\dfrac{5x-1}{6} &= 6\times1 \\
            2(2x+1) - (5x-1) &= 6
        \end{align*}
        \begin{align*}
            4x + 2 - 5x + 1 &= 6
        \end{align*}
        \begin{align*}
            -x + 3 &= 6 \\
            -x &= 3
        \end{align*}
        \begin{align*}
            x &= -3
        \end{align*}

        \item 求父亲现在的年龄。 \\
        $3$ 年前父亲的年龄等于儿子当时年龄的 $7$ 倍，据此列方程。 \\
        设父亲现在 $x$ 岁，则 $3$ 年前父亲 $(x-3)$ 岁，儿子 $(10-3)$ 岁。 \\
        \begin{align*}
            x - 3 &= 7(10-3)
        \end{align*}
        \begin{align*}
            x - 3 &= 49
        \end{align*}
        \begin{align*}
            x &= 52
        \end{align*}
    \end{enumerate}
\end{solution}
%endregion

%region 范围：2.5.2 二元一次方程组及应用题 | 考点：含分数系数的方程组（去分母后用加减消元法）、连等式型方程组 | 难度：中
\begin{question}[points = 10]
    请完成下列各题：

    \begin{enumerate}[label=(\alph*)]
        \item 解方程组：$\begin{cases} \dfrac{y}{3} + \dfrac{x}{2} = 3 \\[2ex] -\dfrac{y}{3}  = 2 - \dfrac{3x}{4}\end{cases}$。\qscore{5}

        \item 解方程：$3x + y = 5x - 2y = 11$。\qscore{5}
    \end{enumerate}
\end{question}

\begin{solution}
    \begin{enumerate}[label=(\alph*)]
        \item 用加减消元法解 $\begin{cases} \dfrac{y}{3} + \dfrac{x}{2} = 3 \\[2ex] -\dfrac{y}{3} + \dfrac{3x}{4} = 2 \end{cases}$。 \\
        先给两式分别去分母，并按 $x$、$y$ 的顺序整理，再把③乘 $2$，使 $y$ 的系数互为相反数，两式相加消去 $y$。 \\
        \begin{align*}
            \text{①}\times6:\quad 3x + 2y &= 18 \quad \text{③} \\
            \text{②}\times12:\quad 9x - 4y &= 24 \quad \text{④}
        \end{align*}
        \hfill (1分) \\
        \begin{align*}
            \text{③}\times2:\quad 6x + 4y &= 36 \quad \text{⑤}
        \end{align*}
        \hfill (1分) \\
        \begin{align*}
            \text{④}+\text{⑤}:\quad 15x &= 60
        \end{align*}
        \hfill (1分) \\
        \begin{align*}
            x &= 4
        \end{align*}
        \hfill (1分) \\
        \begin{align*}
            \text{代入③：}\quad 3\times4 + 2y &= 18 \\
            2y &= 6 \\
            y &= 3
        \end{align*}
        \hfill (1分)

        \item 解连等式方程 $3x + y = 5x - 2y = 11$。 \\
        连等式表示两个代数式都等于 $11$，拆成方程组后用加减消元法求解。 \\
        \begin{align*}
            \begin{cases}
                3x + y = 11 & \text{①} \\
                5x - 2y = 11 & \text{②}
            \end{cases}
        \end{align*}
        \hfill (1分) \\
        \begin{align*}
            \text{①}\times2:\quad 6x + 2y &= 22 \quad \text{③}
        \end{align*}
        \hfill (1分) \\
        \begin{align*}
            \text{②}+\text{③}:\quad 11x &= 33
        \end{align*}
        \hfill (1分) \\
        \begin{align*}
            x &= 3
        \end{align*}
        \hfill (1分) \\
        \begin{align*}
            \text{代入①：}\quad 3\times3 + y &= 11 \\
            y &= 2
        \end{align*}
        \hfill (1分)
    \end{enumerate}
\end{solution}
%endregion
\vfill % 与上方内容保持距离

\noindent
\newlength{\sigwidth}
\settowidth{\sigwidth}{（\hspace{4em}师）}

\begin{tabular*}{\textwidth}{@{\extracolsep{\fill}} l l l l}
    出题： & 审题： & 审题： & 审核： \\[4em]
    \rule{\sigwidth}{0.5pt} &
    \rule{\sigwidth}{0.5pt} &
    \rule{\sigwidth}{0.5pt} &
    \rule{\sigwidth}{0.5pt} \\
    （\hspace{0em} \setterteacher ） &
    （\hspace{0em} \reviewteacher ） &
    （\hspace{0em} \secondreviewteacher ） &
    （\hspace{0em} \auditteacher ） \\
\end{tabular*}

\newpage

\end{document}